% ============================================================
\chapter{Random Fixed Points}
\label{ch:random}
% ============================================================

\begin{intuition}
Random fixed point theory extends classical results to the stochastic setting: one
considers operators $T(\omega, \cdot)$ that depend on a random parameter
$\omega \in \Omega$, and seeks a fixed point $x^*(\omega)$ that is itself measurable in
$\omega$. The main difficulties are of a measurability nature: one must ensure that
constructions (iterations, limits) preserve measurability. These results find applications
in stochastic analysis, mathematical economics, and stochastic game theory.
\end{intuition}

% ------------------------------------------------------------
\section{Random Operators}
% ------------------------------------------------------------

\begin{definition}[Separable measurable space]
\index{separable measurable space}
A topological space $(X, \tau)$ is \emph{separable} if it contains a countable dense
subset. We equip $X$ with its Borel $\sigma$-algebra
$\mathcal{B}(X) = \sigma(\tau)$.
\end{definition}

\begin{definition}[Random operator]
\index{random operator}
Let $(\Omega, \mathcal{A}, \PP)$ be a probability space and $(X, d)$ a separable metric
space. A \emph{random operator} is a map $T\colon \Omega \times X \to X$ such that:
\begin{enumerate}
  \item for every $x \in X$, the map $\omega \mapsto T(\omega, x)$ is
    $(\mathcal{A}, \mathcal{B}(X))$-measurable;
  \item for every $\omega \in \Omega$ (or $\PP$-almost every $\omega$), the map
    $x \mapsto T(\omega, x)$ satisfies certain properties (continuity, contraction, etc.).
\end{enumerate}
\end{definition}

\begin{definition}[Random fixed point]
\index{random fixed point}
A \emph{random fixed point} of $T$ is a random variable $\xi\colon \Omega \to X$
(measurable) such that
\[
  T(\omega, \xi(\omega)) = \xi(\omega) \quad \text{for } \PP\text{-almost every }
  \omega \in \Omega.
\]
\end{definition}

\begin{remark}
The existence of a fixed point for each $\omega$ (by a deterministic theorem) does not
guarantee the existence of a random fixed point: the measurability of
$\omega \mapsto x^*(\omega)$ is a nontrivial problem.
\end{remark}

% ------------------------------------------------------------
\section{Measurable Selection Theorems}
% ------------------------------------------------------------

\begin{definition}[Measurable correspondence]
\index{measurable correspondence}
A correspondence $F\colon \Omega \rightrightarrows X$ is \emph{measurable} if for every
open set $U \subseteq X$, the set $\{\omega : F(\omega) \cap U \neq \emptyset\}$ belongs
to $\mathcal{A}$.
\end{definition}

\begin{theorem}[Kuratowski--Ryll-Nardzewski measurable selection]
\label{thm:measurable-selection}
\index{Kuratowski--Ryll-Nardzewski theorem}
Let $(\Omega, \mathcal{A})$ be a measurable space and $(X, d)$ a complete separable metric
space. If $F\colon \Omega \rightrightarrows X$ is a measurable correspondence with nonempty
closed values, then $F$ admits a \emph{measurable selection}: there exists a measurable
$f\colon \Omega \to X$ with $f(\omega) \in F(\omega)$ for every $\omega$.
\end{theorem}

\begin{proof}[Proof sketch]
One constructs a sequence $(f_n)$ of $\varepsilon_n$-approximate selections (with
$\varepsilon_n \to 0$) by induction. Let $(x_k)_{k \geq 1}$ be a dense sequence in $X$.
For each $n$, partition $\Omega$ into measurable sets
$A_k^n = \{\omega : F(\omega) \cap B(x_k, \varepsilon_n) \neq \emptyset\} \setminus
\bigcup_{j<k} A_j^n$ and set $f_n(\omega) = x_k$ for $\omega \in A_k^n$. The sequence
$(f_n)$ converges uniformly to a measurable selection $f$.
\end{proof}

% ------------------------------------------------------------
\section{Random Banach Fixed Point Theorem}
% ------------------------------------------------------------

\begin{theorem}[Random contraction fixed point]
\label{thm:random-fp-banach}
Let $(\Omega, \mathcal{A}, \PP)$ be a probability space, $(X, d)$ a complete separable
metric space, and $T\colon \Omega \times X \to X$ a random operator such that:
\begin{enumerate}
  \item for every $x \in X$, $\omega \mapsto T(\omega, x)$ is measurable;
  \item for $\PP$-almost every $\omega$, $x \mapsto T(\omega, x)$ is a contraction with
    constant $k(\omega) < 1$.
\end{enumerate}
Then $T$ has a unique random fixed point.
\end{theorem}

\begin{proof}
For each $\omega$, the deterministic Banach theorem provides a unique fixed point
$\xi(\omega) \in X$. It remains to show that $\xi$ is measurable.

Fix $x_0 \in X$ and define the sequence of iterates:
\[
  \xi_0(\omega) = x_0, \qquad \xi_{n+1}(\omega) = T(\omega, \xi_n(\omega)).
\]
By hypothesis, $\xi_0$ is measurable (constant). If $\xi_n$ is measurable, then
$\omega \mapsto T(\omega, \xi_n(\omega))$ is measurable as a composition of measurable
functions (Carath\'eodory's theorem). So $\xi_{n+1}$ is measurable.

For almost every $\omega$, $\xi_n(\omega) \to \xi(\omega)$ (convergence from Banach's
theorem). As a pointwise limit of measurable functions, $\xi$ is measurable.
\end{proof}

\begin{attention}
The separability hypothesis on $X$ is essential to guarantee joint measurability and
measurability of the limit. Without this hypothesis, the result is false in general.
\end{attention}

% ------------------------------------------------------------
\section{Random Schauder Fixed Point Theorem}
% ------------------------------------------------------------

\begin{theorem}[Random Schauder fixed point]
\label{thm:random-fp-schauder}
Let $(\Omega, \mathcal{A})$ be a measurable space, $C$ a closed bounded convex subset of a
separable Banach space $E$, and $T\colon \Omega \times C \to C$ a random operator such that:
\begin{enumerate}
  \item for every $x \in C$, $\omega \mapsto T(\omega, x)$ is measurable;
  \item for every $\omega$, $x \mapsto T(\omega, x)$ is continuous;
  \item for every $\omega$, $T(\omega, \cdot)$ is compact (maps bounded sets to relatively
    compact sets).
\end{enumerate}
Then $T$ has a random fixed point.
\end{theorem}

\begin{proof}[Proof sketch]
By the deterministic Schauder theorem, for each $\omega$, the set of fixed points
$F(\omega) = \{x \in C : T(\omega, x) = x\}$ is nonempty and closed. One shows that the
correspondence $\omega \mapsto F(\omega)$ is measurable. The Kuratowski--Ryll-Nardzewski
measurable selection theorem then provides a measurable selection, which is the desired
random fixed point.
\end{proof}

% ------------------------------------------------------------
\section{Applications in Stochastic Analysis}
% ------------------------------------------------------------

\begin{proposition}[Random differential equations]
Consider the random differential equation
\[
  y'(t, \omega) = f(\omega, t, y(t, \omega)), \qquad y(t_0, \omega) = Y_0(\omega)
\]
where $Y_0\colon \Omega \to \R^n$ is a random variable and $f$ satisfies:
\begin{itemize}
  \item $f(\omega, \cdot, \cdot)$ is continuous in $(t,y)$ for every $\omega$;
  \item $f(\cdot, t, y)$ is measurable in $\omega$ for all $t, y$;
  \item $f$ is Lipschitz in $y$ uniformly in $\omega$.
\end{itemize}
Then there exists a unique solution $y(\cdot, \omega)$ that is a measurable process.
\end{proposition}

\begin{proof}
The Picard operator
$T(\omega, u)(t) = Y_0(\omega) + \int_{t_0}^t f(\omega, s, u(s)) ds$ is a contractive
random operator on $C([t_0-\delta, t_0+\delta], \R^n)$ equipped with the Bielecki norm.
Theorem~\ref{thm:random-fp-banach} applies.
\end{proof}

\begin{proposition}[Stochastic integral equations]
The random Fredholm integral equation
\[
  u(\omega, t) = f(\omega, t) + \lambda \int_a^b K(\omega, t, s) u(\omega, s) \, ds
\]
has a unique random fixed point whenever
$\abs{\lambda} < 1/((b-a) \sup_\omega \norm{K(\omega, \cdot, \cdot)}_\infty)$.
\end{proposition}

% ------------------------------------------------------------
\section{Itoh's Theorem}
% ------------------------------------------------------------

\begin{theorem}[Itoh]
\label{thm:itoh}
\index{Itoh's theorem}
Let $(\Omega, \mathcal{A})$ be a complete measurable space, $C$ a closed separable subset
of a Banach space, and $T\colon \Omega \times C \to C$ a nonexpansive random operator
($d(T(\omega,x), T(\omega,y)) \leq d(x,y)$) such that for every $\omega$, $T(\omega, C)$
is contained in a compact subset of $C$. Then $T$ has a random fixed point.
\end{theorem}

\begin{keyformulas}[title=\textbf{Summary of random fixed point theorems}]
\begin{center}
\renewcommand{\arraystretch}{1.3}
\begin{tabular}{|l|l|l|}
\hline
\textbf{Theorem} & \textbf{Hypothesis on $T(\omega, \cdot)$} & \textbf{Uniqueness} \\
\hline
Random Banach & Contraction & Yes \\
Random Schauder & Continuous, compact & No \\
Itoh & Nonexpansive, compact image & No \\
\hline
\end{tabular}
\end{center}
\end{keyformulas}

% ------------------------------------------------------------
\section{Exercises}
% ------------------------------------------------------------

\begin{exercise}
Let $\Omega = [0,1]$ with the Lebesgue measure and $X = \R$. Define
$T(\omega, x) = \omega x / 2$. Verify that $T$ is a contractive random operator and
determine the random fixed point explicitly.
\end{exercise}

\begin{exercise}
Let $T(\omega, x) = a(\omega) x + b(\omega)$ where $a, b\colon \Omega \to \R$ are
measurable with $\abs{a(\omega)} < 1$ a.s. Show that the random fixed point is
$\xi(\omega) = b(\omega) / (1 - a(\omega))$. Verify its measurability.
\end{exercise}

\begin{exercise}
Let $(X, d)$ be a compact metric space and $T\colon \Omega \times X \to X$ a continuous
random operator. Show that the correspondence $F(\omega) = \{x : T(\omega, x) = x\}$ is
measurable.
\emph{Hint: show that $\{\omega : F(\omega) \cap U \neq \emptyset\}$ is measurable for
every open set $U$.}
\end{exercise}

\begin{exercise}
Consider the random ODE $y' = A(\omega) y$, $y(0) = 1$, where $A(\omega)$ is a real
random variable. Reformulate as a fixed point problem and apply the random contraction
fixed point theorem.
\end{exercise}

\begin{exercise}
Give an example of a random operator $T\colon \Omega \times [0,1] \to [0,1]$ continuous
in $x$ and measurable in $\omega$, which has a fixed point for each $\omega$, but for
which no measurable selection of fixed points exists.
\emph{Hint: use a non-complete measurable space $(\Omega, \mathcal{A})$.}
\end{exercise}

\begin{exercise}[Random iterations]
\index{iterated function system}
Let $(T_1, \ldots, T_k)$ be a family of contractions on $\R^n$ with probabilities
$(p_1, \ldots, p_k)$. Define $x_{n+1} = T_{\sigma_n}(x_n)$ where $(\sigma_n)$ is i.i.d.\
with law $(p_i)$. Show that the sequence $(x_n)$ converges a.s.\ to an invariant measure
(random attractor), independently of the initial point $x_0$.
\end{exercise}
